\selectlanguage{spanish}

\renewcommand{\articuloidiomaxml}{es}
\renewcommand{\articulotitulo}{Las deudas pendientes de la digitalización en América Latina}
\renewcommand{\articuloautores}{Maria Mora}
\renewcommand{\articulodoi}{10.31637/epsir-2023-898}
\renewcommand{\articulotipo}{EDITORIAL}
\renewcommand{\articulorevista}{Revista Latinoamericana de Estudios Editoriales y Cultura Escrita}
\renewcommand{\articuloissne}{2591-3905}
\renewcommand{\articulovolumen}{01}
\renewcommand{\articulonumero}{01}
\renewcommand{\articulofecha}{ 2026}
\renewcommand{\articulopagina}{1}

\setcounter{page}{1}
\thispagestyle{firstpage}

% --- TÍTULO A ANCHO COMPLETO ---
\noindent\begin{minipage}[t]{\dimexpr\textwidth+\marginparsep+\marginparwidth\relax}%
{\sffamily\bfseries\Large\articulotitulo\par}%
\end{minipage}\par
\vspace{3mm}%
\renewcommand{\articulotitulotrans}{The outstanding debts of digitalization in Latin America}
\noindent\begin{minipage}[t]{\dimexpr\textwidth+\marginparsep+\marginparwidth\relax}%
{\sffamily\itshape\normalsize\articulotitulotrans\par}%
\end{minipage}\par
\vspace{4mm}%

% --- BLOQUE LATERAL DE METADATOS ---
\begin{textblock*}{68mm}(\gbbloquex,92mm)
{\setlength{\leftskip}{0pt}\setlength{\parindent}{0pt}%
\noindent\begin{minipage}[t]{68mm}
\sffamily\small\setlength{\parskip}{1pt}
{\bfseries\color{azulrevista}Maria Mora}\par
{\scriptsize\texttt{\href{https://orcid.org/0000-0003-0012-4613}{https://orcid.org/0000-0003-0012-4613}}}\par
departamento, Facultad de Filosofía y Letras, Buenos Aires, Argentina\par
\vspace{6pt}
{\bfseries\color{azulrevista!70}Derechos de autor\'ia:}\par
\textcopyright\ 2026 Mora, M. Publicado por Revista Latinoamericana de Estudios Editoriales y Cultura Escrita. Se permite el uso, distribución y reproducción sin fines comerciales, siempre que se cite la obra original. (\href{https://creativecommons.org/licenses/by-nc/4.0/}{https://creativecommons.org/licenses/by-nc/4.0/})\par\vspace{6pt}
{\bfseries\color{azulrevista!70}Publicación:} 2026, vol.~01, n\'um.~01, pp.~1--2\par\vspace{6pt}
{\bfseries\color{azulrevista!70}URL:} {\scriptsize\texttt{\href{https://redeseditoriales.org/revista-latinoamericana-de-estudios-editoriales/}{https://redeseditoriales.org/revista-latinoamericana-de-estudios-editoriales/}}}\par\vspace{6pt}
{\bfseries\color{azulrevista!70}DOI:} {\scriptsize\texttt{\href{https://doi.org/10.31637/epsir-2023-898}{10.31637/epsir-2023-898}}}\par\vspace{6pt}
{\bfseries\color{azulrevista!70}Licencia:} Creative Commons Attribution 4.0 International (CC BY 4.0)\par
{\scriptsize\texttt{\href{https://creativecommons.org/licenses/by-nc/4.0/}{https://creativecommons.org/licenses/by-nc/4.0/}}}\par
\vspace{6pt}
{\bfseries\color{azulrevista!70}Fechas}\par\vspace{2pt}
Online: 26/06/2026\par
\vspace{6pt}
\vspace{6pt}
{\bfseries\color{azulrevista!70}Compartir:}\par\vspace{3pt}
{\large
\href{https://bsky.app/intent/compose?text=https%3A%2F%2Fredeseditoriales.org%2Frevista-latinoamericana-de-estudios-editoriales%2F}{\includesvg[height=0.9em]{/home/maria/.gbpublisher/fonts/bluesky}}~
\href{mailto:?subject=Las%20deudas%20pendientes%20de%20la%20digitalizaci%C3%B3n%20en%20Am%C3%A9rica%20Latina&body=https%3A%2F%2Fredeseditoriales.org%2Frevista-latinoamericana-de-estudios-editoriales%2F}{\includesvg[height=0.9em]{/home/maria/.gbpublisher/fonts/envelope}}~
\href{https://www.facebook.com/sharer/sharer.php?u=https%3A%2F%2Fredeseditoriales.org%2Frevista-latinoamericana-de-estudios-editoriales%2F}{\includesvg[height=0.9em]{/home/maria/.gbpublisher/fonts/facebook}}~
\href{https://www.linkedin.com/sharing/share-offsite/?url=https%3A%2F%2Fredeseditoriales.org%2Frevista-latinoamericana-de-estudios-editoriales%2F}{\includesvg[height=0.9em]{/home/maria/.gbpublisher/fonts/linkedin}}~
\href{https://www.reddit.com/submit?url=https%3A%2F%2Fredeseditoriales.org%2Frevista-latinoamericana-de-estudios-editoriales%2F&title=Las%20deudas%20pendientes%20de%20la%20digitalizaci%C3%B3n%20en%20Am%C3%A9rica%20Latina}{\includesvg[height=0.9em]{/home/maria/.gbpublisher/fonts/reddit}}~
\href{https://www.researchgate.net/search?q=Las%20deudas%20pendientes%20de%20la%20digitalizaci%C3%B3n%20en%20Am%C3%A9rica%20Latina}{\includesvg[height=0.9em]{/home/maria/.gbpublisher/fonts/researchgate}}~
\href{https://t.me/share/url?url=https%3A%2F%2Fredeseditoriales.org%2Frevista-latinoamericana-de-estudios-editoriales%2F&text=Las%20deudas%20pendientes%20de%20la%20digitalizaci%C3%B3n%20en%20Am%C3%A9rica%20Latina}{\includesvg[height=0.9em]{/home/maria/.gbpublisher/fonts/telegram}}~
\href{https://wa.me/?text=https%3A%2F%2Fredeseditoriales.org%2Frevista-latinoamericana-de-estudios-editoriales%2F}{\includesvg[height=0.9em]{/home/maria/.gbpublisher/fonts/whatsapp}}~
\href{https://twitter.com/intent/tweet?url=https%3A%2F%2Fredeseditoriales.org%2Frevista-latinoamericana-de-estudios-editoriales%2F&text=Las%20deudas%20pendientes%20de%20la%20digitalizaci%C3%B3n%20en%20Am%C3%A9rica%20Latina}{\includesvg[height=0.9em]{/home/maria/.gbpublisher/fonts/x-twitter}}
}\par
\end{minipage}%
}% FIN GRUPO leftskip
\end{textblock*}


\begin{tcolorbox}[colback=brown!8!white,colframe=brown!8!white,boxrule=0pt,arc=2pt,left=4pt,right=4pt,top=3pt,bottom=3pt,before skip=4pt,after skip=4pt]
\begin{otherlanguage}{spanish}
\sffamily\small \textbf{Resumen:} La pandemia aceleró la digitalización en América Latina en áreas como la educación, los trámites públicos y la salud. Sin embargo, este proceso evidenció y profundizó las fracturas de una región con infraestructura y competencias marcadamente desiguales entre las grandes ciudades y las zonas rurales o periféricas.
Aunque recientemente aumentó la penetración general de internet, persiste una "brecha digital de segundo nivel". Los sectores vulnerables acceden mayoritariamente mediante teléfonos móviles con planes de datos limitados, lo que restringe severamente su desarrollo educativo, laboral y ciudadano. De hecho, la conectividad de banda ancha resulta económicamente inviable para los más pobres, llegando a absorber hasta un 14\% de sus ingresos mensuales.
Proveer únicamente infraestructura tecnológica es insuficiente. Una verdadera inclusión digital requiere políticas públicas integrales que atiendan la conectividad, la formación y los contenidos, contemplando las desigualdades de ingreso, género, etnia y territorio. La tecnología, por sí sola, no resolverá las deudas históricas de la región.
\par\smallskip\noindent\textbf{Palabras clave:} Brecha digital de segundo nivel, Digitalización, América Latina, Políticas públicas, Infraestructura digital, Desigualdad social
\end{otherlanguage}
\end{tcolorbox}

\begin{tcolorbox}[colback=brown!8!white,colframe=brown!8!white,boxrule=0pt,arc=2pt,left=4pt,right=4pt,top=3pt,bottom=3pt,before skip=4pt,after skip=4pt]
\begin{otherlanguage}{english}
\sffamily\small \textbf{Abstract:} The pandemic accelerated digitalization in Latin America across areas such as education, public procedures, and healthcare. However, this process exposed and deepened the fractures of a region with markedly unequal infrastructure and skills between major cities and rural or peripheral areas.
Although general internet penetration has recently increased, a "second-level digital divide" persists. Vulnerable sectors access the internet primarily through mobile phones with limited data plans, severely restricting their educational, professional, and civic development. In fact, broadband connectivity remains financially unviable for the poorest, consuming up to 14\% of their monthly income.
Providing technological infrastructure alone is insufficient. True digital inclusion requires comprehensive public policies that address connectivity, training, and content, while accounting for pre-existing inequalities of income, gender, ethnicity, and territory. Technology, on its own, will not resolve the region's historical debts.
\par\smallskip\noindent\textbf{Keywords:} Second-level digital divide, Digitalization, Latin America, Public policies, Digital infrastructure, Social inequality
\end{otherlanguage}
\end{tcolorbox}

\bigskip
\noindent\rule{\linewidth}{0.4pt}
\bigskip

\section{Introducción}
La pandemia de COVID-19 funcionó como un acelerador involuntario de procesos que llevaban años en gestación. En cuestión de semanas, sistemas educativos enteros migraron a plataformas virtuales, trámites administrativos que exigían presencialidad se trasladaron a portales web y la telemedicina dejó de ser una promesa para convertirse en una necesidad imperiosa. Sin embargo, esa aceleración no hizo más que volver más visibles las fracturas que ya existían en el tejido social de la región.

América Latina llegó a la pandemia con una infraestructura digital profundamente desigual. Mientras las capitales y las grandes áreas metropolitanas disponían de conexiones de banda ancha razonablemente estables, vastas zonas rurales e incluso periferias urbanas carecían de conectividad básica. La brecha no era solo de acceso físico: se manifestaba también en las competencias digitales de la población, en la disponibilidad de dispositivos adecuados y en la capacidad institucional para sostener servicios en línea de manera continua y segura.

Tres años después del momento más agudo de la crisis sanitaria, el balance resulta ambivalente. Por un lado, se registró una expansión significativa de la conectividad: varios países de la región superaron tasas de penetración de Internet que parecían lejanas a comienzos de 2020. Por otro, las desigualdades de uso —lo que la literatura especializada denomina “brecha digital de segundo nivel”— no solo persistieron sino que, en algunos casos, se profundizaron. Los sectores de menores ingresos accedieron mayoritariamente a través de teléfonos móviles con planes de datos limitados, lo que restringía severamente sus posibilidades de participar en actividades educativas, laborales o de ejercicio ciudadano que demandan conexiones estables y pantallas de mayor tamaño.\footnote{Según estimaciones de la CEPAL, en 2022 el costo del servicio de
  banda ancha para la población del primer quintil de ingresos
  representaba en promedio el 14\% de su ingreso mensual en la región, una
  proporción que vuelve inviable cualquier pretensión de universalidad
  efectiva.}

El problema, conviene insistir, no es exclusivamente tecnológico. La digitalización que no va acompañada de políticas públicas integrales —que atiendan simultáneamente la conectividad, la formación, el diseño institucional y la producción de contenidos pertinentes— corre el riesgo de reproducir y amplificar las desigualdades preexistentes. En una región donde la matriz de la desigualdad opera simultáneamente por ingreso, género, etnia, territorio y edad, suponer que la mera provisión de infraestructura resuelve el acceso es una simplificación que ningún responsable de política pública debería permitirse.

Este número de la revista reúne contribuciones que abordan distintas aristas de esta problemática. Los trabajos que aquí se presentan comparten una premisa: la transformación digital de nuestras sociedades no puede analizarse al margen de las condiciones materiales y simbólicas en las que se despliega. Esperamos que su lectura contribuya a enriquecer un debate que, lejos de cerrarse, se vuelve más urgente con cada innovación tecnológica que promete —una vez más— resolver por sí sola las deudas históricas de la región.


